\section{Structure et organigramme du LaBRI}
\label{annexe:labri}

Le Laboratoire Bordelais de Recherche en Informatique (LaBRI), UMR 5800, au sein duquel s'est déroulé ce stage de fin d'études au sein de l'équipe BKB, est structuré autour de plusieurs grands pôles :
\begin{itemize}
    \item \textbf{Direction :} Pascal Desbarats (Directeur par intérim), Nicolas Hanusse (Directeur adjoint par intérim), Serge Chaumette (Chargé de mission transfert et valorisation).
    \item \textbf{Administration} (Administratrice d'unité : Christine Richard) :
    \begin{itemize}
        \item \textit{Équipe administrative} (resp. Cathy Roubineau) : Accueil (Claire Douvier, Laurence Pinosa), Assistante administrative (Safia Amsili), Assistante de communication (Auriane Dantès).
        \item \textit{Équipe financière} (resp. Chloé Godard) : Gestionnaires (Isabelle Garcia, Anthony Gau, Emmanuelle Lesage, Elia Meyre).
        \item \textit{Assistante du SCRIME} : Annick Mersier.
    \end{itemize}
    \item \textbf{Recherche} (structurée en plusieurs départements) :
    \begin{itemize}
        \item \textit{Image et Son} (Fabien Baldacci)
        \item \textit{Combinatoire et Algorithmique} (Cyril Gavoille)
        \item \textit{Systèmes et Données} (Jean-Rémy Falleri), département auquel appartient l'équipe BKB.
        \item \textit{Méthodes et Modèles formels} (Laurent Bienvenu)
        \item \textit{Supports et AlgoriThmes pour les Applications Numériques hAutes performanceS} (Pierre Ramet)
    \end{itemize}
    \item \textbf{Appui à la recherche :}
    \begin{itemize}
        \item \textit{Ingénieurs affectés sur projet} : Nathalie Furmento, Frédéric Lalanne, Boris Mansencal, Gérald Point.
        \item \textit{Équipe soutien "Systèmes - Réseaux"} (resp. Pascal Ung) : Ingénieurs ASR (Michaël Fontaine, Emmanuel Fontaine, Gabriel Larrouy), Technicien ASR (Alexandre Issouli), Ingénieur DEV-Web (Pierre Lacroix).
    \end{itemize}
    \item \textbf{Autre :} Assistant de prévention (Gérald Point).
\end{itemize}

L'organigramme structurel simplifié correspondant est présenté dans la Figure~\ref{fig:organigramme_labri}.

\begin{figure}[H]
  \centering
  \begin{tikzpicture}[
    box/.style={draw=ubblue, fill=softgray, thick, rounded corners, align=center, text width=3.8cm, font=\small, minimum height=1cm},
    dirbox/.style={box, fill=ubblue!10, text width=6cm},
    deptbox/.style={box, draw=green!60!black, fill=green!5, text width=2.45cm, font=\scriptsize},
    adminbox/.style={box, draw=orange!70!black, fill=orange!5, text width=3.5cm, font=\scriptsize},
    supportbox/.style={box, draw=purple!70!black, fill=purple!5, text width=3.5cm, font=\scriptsize},
    line/.style={draw, thick, -latex}
  ]

  % Nodes
  \node (dir) [dirbox] {\textbf{Direction du LaBRI} \\ \scriptsize P. Desbarats (Dir. int.), N. Hanusse (Dir. adj. int.), S. Chaumette (Transfert)};

  \node (admin) [adminbox, below left=1.2cm and 0.5cm of dir] {\textbf{Administration} \\ (Christine Richard) \\ \vspace{0.1cm} \raggedright \scriptsize $\bullet$ Éq. Administrative \\ (C. Roubineau) \\ $\bullet$ Éq. Financière \\ (C. Godard) \\ $\bullet$ SCRIME \\ (A. Mersier)};

  \node (support) [supportbox, below right=1.2cm and 0.5cm of dir] {\textbf{Appui à la Recherche} \\ \vspace{0.1cm} \raggedright \scriptsize $\bullet$ Systèmes \& Réseaux \\ (P. Ung) \\ $\bullet$ Ingénieurs projets \\ (N. Furmento, F. Lalanne, B. Mansencal, G. Point) \\ $\bullet$ Assistant Prévention \\ (G. Point)};

  \node (coord) [below=5.1cm of dir, inner sep=0pt, minimum size=0pt] {};

  % Departements de recherche
  \node (dept3) [deptbox, below=0.6cm of coord] {\textbf{Systèmes \& Données} \\ (J.-R. Falleri) \\ \vspace{0.05cm} \textit{Équipe BKB (accueil)}};
  \node (dept2) [deptbox, left=0.2cm of dept3] {\textbf{Combinatoire \& Algorithmique} \\ (C. Gavoille)};
  \node (dept1) [deptbox, left=0.2cm of dept2] {\textbf{Image \& Son} \\ (F. Baldacci)};
  \node (dept4) [deptbox, right=0.2cm of dept3] {\textbf{Méthodes \& Modèles} \\ (L. Bienvenu)};
  \node (dept5) [deptbox, right=0.2cm of dept4] {\textbf{Satanas} \\ (P. Ramet)};

  % Lines
  \path [draw, thick] (dir) -- (coord);
  \path [draw, thick] (dir) -| (admin);
  \path [draw, thick] (dir) -| (support);
  \path [draw, thick] (coord) -| (dept1);
  \path [draw, thick] (coord) -| (dept2);
  \path [draw, thick] (coord) -- (dept3);
  \path [draw, thick] (coord) -| (dept4);
  \path [draw, thick] (coord) -| (dept5);

  \end{tikzpicture}
  \caption{Organigramme structurel simplifié du LaBRI.}
  \label{fig:organigramme_labri}
\end{figure}

\newpage
\section{Détails des datasets de l'état de l'art}
\label{annexe:datasets}

Cette section fournit les détails des huit jeux de données scRNA-seq du sous-ensemble commun (``common-8'') utilisés dans l'étude, y compris l'espèce, les tailles (cellules et gènes), le nombre de lots (batchs), le nombre de populations cellulaires rares, ainsi que leur niveau d'effet de lot (mesuré par PCR sur PCA) et de déséquilibre de classe (mesuré par l'indice de Gini).

\begin{table}[H]
\centering
\footnotesize
\setlength{\tabcolsep}{4pt}
\begin{tabular}{l l c c c c c c}
\toprule
\textbf{Dataset} & \textbf{Espèce} & \textbf{Cellules} & \textbf{Gènes} & \textbf{Lots} & \textbf{Rares ($<5\%$)} & \textbf{PCR/PCA} & \textbf{Gini} \\
\midrule
BBAG094 Zeisel & Souris & 3 006 & 19 972 & 10 & 4 & 0,499 & 0,543 \\
BBAG094 spleen & Souris & 9 552 & 23 341 & 2 & 3 & 0,022 & 0,644 \\
Baron Human Pancreas & Humain & 8 569 & 20 125 & 4 & 9 & 0,048 & 0,644 \\
Human testis GSE112013 & Humain & 6 490 & 27 477 & 6 & 4 & 0,025 & 0,304 \\
Kang PBMC & Humain & 24 679 & 35 635 & 16 & 1 & 0,069 & 0,500 \\
Macaque retina bipolar & Macaque & 30 302 & 36 162 & 30 & 4 & 0,184 & 0,388 \\
Paul15 bone marrow & Souris & 2 730 & 3 451 & 1 & 9 & 0,000 & 0,410 \\
Tabula Muris liver & Souris & 2 859 & 22 966 & 1 & 5 & 0,000 & 0,578 \\
\bottomrule
\end{tabular}
\caption{Récapitulatif des huit jeux de données du sous-ensemble commun (common-8) utilisés dans nos expériences.}
\label{tab:datasets}
\end{table}

\subsection*{Description des métriques de caractérisation des jeux de données}

\begin{itemize}
    \item \textbf{PCR sur PCA (Principal Component Regression on PCA) :} Cette métrique quantifie l'importance globale de l'effet de lot dans un jeu de données en mesurant la part de variance expliquée par la variable de lot (batch) dans l'espace des composantes principales. Une analyse en composantes principales (PCA) est d'abord réalisée sur la matrice d'expression. Pour chaque composante principale, une régression linéaire est ajustée en utilisant le lot comme variable explicative afin d'obtenir un coefficient de détermination $R^2$. La métrique finale est la somme de ces $R^2$ pondérée par la variance expliquée par chaque composante. Un score proche de 0 indique l'absence d'effet de lot, tandis qu'un score proche de 1 indique que l'effet de lot domine la variance des données.
    \item \textbf{Indice de Gini des classes :} Cette métrique mesure le déséquilibre de la distribution des tailles des populations cellulaires (les classes de la vérité terrain) dans le jeu de données. Issu de la théorie économique pour mesurer les inégalités de richesse, cet indice est adapté ici pour mesurer la dispersion statistique des effectifs des différents types cellulaires. Un indice de 0 correspond à une répartition parfaitement équilibrée (toutes les classes ont le même nombre de cellules), tandis qu'un indice proche de 1 traduit un déséquilibre extrême (une classe ultra-majoritaire face à de nombreuses classes très minoritaires).
\end{itemize}

\newpage
\section{Tableau des métriques d'évaluation}
\label{annexe:metriques}

Cette section regroupe les définitions, formules mathématiques et interprétations des métriques utilisées pour évaluer et comparer les performances des algorithmes de clustering, tant sur l'aspect global que sur la détection des classes rares et ultra-rares.

\begin{table}[H]
    \centering
    \scriptsize
    \renewcommand{\arraystretch}{1.2}
    \begin{tabularx}{\textwidth}{@{}p{0.15\textwidth}p{0.17\textwidth}p{0.35\textwidth}X@{}}
        \toprule
        \textbf{Métrique} & \textbf{Rôle} & \textbf{Formule ou principe} & \textbf{Interprétation} \\
        \midrule
        ARI & Clustering global & \(ARI=\frac{RI-\mathrm{E}[RI]}{\max(RI)-\mathrm{E}[RI]}\), où \(RI\) mesure l'accord entre paires de cellules. & Mesure si deux cellules regroupées ensemble dans la vérité biologique le sont aussi dans la prédiction, en corrigeant l'accord attendu au hasard. \\
        NMI & Clustering global & \(NMI(Y,\hat{Y})=\frac{2I(Y;\hat{Y})}{H(Y)+H(\hat{Y})}\) & Mesure la quantité d'information partagée entre les annotations biologiques et les groupes prédits. \\
        ACC & Classification après appariement des clusters & \(ACC=\frac{1}{n}\max_{\pi}\sum_{i=1}^{n}\mathbf{1}\{y_i=\pi(\hat{y}_i)\}\) & Donne la proportion globale de cellules correctement attribuées après correspondance entre clusters et classes. \\
        BalancedACC & Classification équilibrée & BalancedACC\({}=\frac{1}{|\mathcal{C}|}\sum_{c\in\mathcal{C}}\frac{TP_c}{N_c}\) & Calcule la moyenne des rappels par classe, afin qu'une classe très abondante ne masque pas les erreurs sur les classes petites. \\
        RareACC & Classification des classes rares & RareACC\({}=\frac{1}{|\mathcal{R}|}\sum_{c\in\mathcal{R}}\frac{TP_c}{N_c}\), avec \(\mathcal{R}=\{c\,|\,N_c/n<5\,\%\}\) & Évalue spécifiquement les types cellulaires rares, qui représentent moins de 5\,\% des cellules. \\
        UltraRareACC & Classification des classes ultra-rares & UltraRareACC\({}=\frac{1}{|\mathcal{U}|}\sum_{c\in\mathcal{U}}\frac{TP_c}{N_c}\), avec \(\mathcal{U}=\{c\,|\,N_c/n<1\,\%\}\) & Mesure la capacité de la méthode à préserver les populations les plus minoritaires. \\
        Temps & Coût computationnel & \(t_{\mathrm{exec}}=t_{\mathrm{fin}}-t_{\mathrm{debut}}\) & Permet de comparer les méthodes non seulement sur leur qualité de clustering, mais aussi sur leur coût pratique. \\
        \bottomrule
    \end{tabularx}
    \caption{Métriques utilisées pour comparer les méthodes. \(TP_c\) correspond au nombre de cellules de la classe \(c\) correctement attribuées après appariement, et \(N_c\) au nombre total de cellules de cette classe.}
    \label{tab:metriques_evaluation}
\end{table}

Dans les formules ci-dessus, \(y_i\) désigne l'annotation biologique de la cellule \(i\), \(\hat{y}_i\) le cluster prédit par la méthode, \(n\) le nombre total de cellules, \(\mathcal{C}\) l'ensemble des classes biologiques et \(\pi\) l'appariement optimal obtenu par l'algorithme hongrois.

\newpage
\section{Résultats complets sur les huit jeux de données communs}
\label{annexe:resultats_common8}

Cette annexe présente les résultats complets utilisés pour construire la figure principale de la section scRAW. Les huit jeux de données correspondent au sous-ensemble commun utilisé pendant l'évaluation généraliste : BBAG094 Zeisel, BBAG094 spleen, Baron Human Pancreas, Human testis GSE112013, Kang PBMC, Macaque retina bipolar, Paul15 bone marrow et Tabula Muris liver. La première figure reprend tous les algorithmes de chaque famille principale, avec \textbf{scRAW} ajouté comme repère dans chaque colonne ; la seconde ajoute les expériences complémentaires où Harmony est appliqué en correction de batch.

\begin{figure}[H]
    \centering
    \includegraphics[width=\textwidth,height=0.93\textheight,keepaspectratio]{scraw_common8_family_panel/common8_family_all_methods_panel.pdf}
    \caption{Résultats complets par famille sur les huit jeux de données communs. Les boîtes représentent la distribution des scores disponibles par méthode sur les jeux de données ; pour l'UltraRareACC, les jeux sans classe ultra-rare ne contribuent pas à la distribution. Les méthodes sont ordonnées par moyenne croissante dans chaque sous-panneau. \textbf{scRAW} est ajouté comme repère dans chaque famille.}
    \label{fig:common8_family_all_methods}
\end{figure}
Les tableaux suivants présentent les résultats bruts détaillés par jeu de données pour chaque métrique d'évaluation pour les méthodes principales :
\begin{table}[H]
  \centering
  \fontsize{5.5pt}{7.0pt}\selectfont
  \setlength{\tabcolsep}{1.0pt}
  \resizebox{\textwidth}{!}{%
  \begin{tabular}{l|c|ccccc|cccc|cccc}
    \toprule
    \textbf{Jeu de données} & \multicolumn{1}{c|}{\textbf{scRAW}} & \multicolumn{5}{c|}{\textbf{Rare Specific}} & \multicolumn{4}{c|}{\textbf{Généralistes}} & \multicolumn{4}{c}{\textbf{Correction Batch}} \\
    & \textbf{scRAW} & scAIDE & CellSIUS & DeepScena & scCAD & GiniClust & scVI & scMAE & PCA+Leiden & scNAME & Harmony & ComBat & DESC & Scanorama \\
    \midrule
    Zeisel & 0,430 & \underline{0,696} & 0,572 & 0,316 & 0,358 & 0,036 & 0,513 & 0,692 & 0,632 & 0,645 & 0,388 & 0,000 & \textbf{0,752} & 0,632 \\
    Spleen & 0,870 & \textbf{0,934} & 0,465 & 0,517 & 0,416 & 0,141 & 0,572 & \underline{0,920} & 0,433 & 0,689 & 0,518 & 0,580 & 0,336 & 0,448 \\
    Baron Pancreas & \underline{0,916} & 0,794 & 0,666 & 0,572 & 0,605 & 0,052 & 0,621 & 0,770 & 0,666 & 0,708 & 0,851 & 0,653 & \textbf{0,918} & 0,621 \\
    Human Testis & 0,611 & 0,611 & 0,581 & 0,492 & 0,399 & 0,043 & 0,594 & 0,651 & 0,626 & 0,631 & \underline{0,652} & 0,569 & \textbf{0,668} & 0,471 \\
    Kang PBMC & 0,546 & \textbf{0,811} & 0,775 & 0,773 & 0,428 & 0,052 & \underline{0,805} & 0,779 & 0,774 & 0,751 & 0,774 & 0,773 & 0,533 & 0,782 \\
    Macaque Retina & 0,924 & 0,873 & 0,497 & 0,046 & 0,425 & 0,023 & \underline{0,965} & 0,801 & 0,497 & 0,556 & 0,941 & \textbf{0,971} & 0,805 & 0,351 \\
    Bone Marrow & 0,283 & 0,328 & 0,341 & 0,277 & 0,224 & 0,007 & 0,283 & \textbf{0,386} & 0,311 & \underline{0,359} & 0,312 & 0,313 & 0,282 & 0,311 \\
    TM Liver & 0,683 & 0,627 & 0,365 & \textbf{0,824} & 0,443 & 0,103 & 0,519 & 0,491 & 0,554 & 0,589 & 0,535 & 0,554 & \underline{0,700} & 0,554 \\
    \midrule
    \textbf{Moyenne} & 0,658 & \textbf{0,709} & 0,533 & 0,477 & 0,412 & 0,057 & 0,609 & \underline{0,686} & 0,562 & 0,616 & 0,621 & 0,552 & 0,624 & 0,521 \\
    \bottomrule
  \end{tabular}
  }
  \caption[Scores ARI bruts (méthodes principales)]{Résultats bruts de Rand Index Ajusté (ARI) pour les méthodes principales sur les 8 jeux de données communs. La meilleure performance par jeu de données est en gras, et la deuxième est soulignée.}
  \label{tab:common8_primary_raw_ARI}
\end{table}

\begin{table}[H]
  \centering
  \fontsize{5.5pt}{7.0pt}\selectfont
  \setlength{\tabcolsep}{1.0pt}
  \resizebox{\textwidth}{!}{%
  \begin{tabular}{l|c|ccccc|cccc|cccc}
    \toprule
    \textbf{Jeu de données} & \multicolumn{1}{c|}{\textbf{scRAW}} & \multicolumn{5}{c|}{\textbf{Rare Specific}} & \multicolumn{4}{c|}{\textbf{Généralistes}} & \multicolumn{4}{c}{\textbf{Correction Batch}} \\
    & \textbf{scRAW} & scAIDE & CellSIUS & DeepScena & scCAD & GiniClust & scVI & scMAE & PCA+Leiden & scNAME & Harmony & ComBat & DESC & Scanorama \\
    \midrule
    Zeisel & 0,641 & \textbf{0,835} & 0,625 & 0,552 & 0,343 & 0,297 & 0,716 & 0,832 & 0,790 & 0,779 & 0,606 & 0,315 & \underline{0,834} & 0,790 \\
    Spleen & 0,829 & \underline{0,957} & 0,687 & 0,729 & 0,557 & 0,673 & 0,796 & \textbf{0,959} & 0,606 & 0,869 & 0,687 & 0,733 & 0,518 & 0,641 \\
    Baron Pancreas & \textbf{0,934} & 0,750 & 0,731 & 0,660 & 0,537 & 0,310 & 0,712 & 0,731 & 0,727 & 0,739 & 0,891 & 0,722 & \underline{0,929} & 0,687 \\
    Human Testis & 0,718 & 0,726 & 0,664 & 0,596 & 0,315 & 0,217 & 0,725 & 0,720 & \underline{0,744} & 0,709 & 0,730 & 0,642 & \textbf{0,820} & 0,562 \\
    Kang PBMC & 0,706 & 0,732 & 0,750 & 0,749 & 0,408 & 0,421 & \textbf{0,814} & 0,744 & \underline{0,754} & 0,747 & 0,753 & 0,753 & 0,643 & 0,744 \\
    Macaque Retina & 0,937 & 0,900 & 0,524 & 0,273 & 0,560 & 0,220 & \underline{0,973} & 0,813 & 0,525 & 0,601 & 0,965 & \textbf{0,985} & 0,806 & 0,399 \\
    Bone Marrow & 0,416 & 0,466 & 0,488 & 0,316 & 0,287 & 0,153 & 0,403 & \textbf{0,511} & 0,473 & \underline{0,492} & 0,472 & 0,468 & 0,359 & 0,473 \\
    TM Liver & 0,746 & 0,696 & 0,497 & \textbf{0,822} & 0,401 & 0,387 & 0,688 & 0,603 & 0,719 & 0,706 & 0,709 & 0,719 & \underline{0,758} & 0,719 \\
    \midrule
    \textbf{Moyenne} & \underline{0,741} & \textbf{0,758} & 0,621 & 0,587 & 0,426 & 0,335 & 0,728 & 0,739 & 0,667 & 0,705 & 0,727 & 0,667 & 0,708 & 0,627 \\
    \bottomrule
  \end{tabular}
  }
  \caption{Résultats bruts de Précision Globale (ACC) pour les méthodes principales sur les 8 jeux de données communs. La meilleure performance par jeu de données est en gras, et la deuxième est soulignée.}
  \label{tab:common8_primary_raw_ACC}
\end{table}

\begin{table}[H]
  \centering
  \fontsize{5.5pt}{7.0pt}\selectfont
  \setlength{\tabcolsep}{1.0pt}
  \resizebox{\textwidth}{!}{%
  \begin{tabular}{l|c|ccccc|cccc|cccc}
    \toprule
    \textbf{Jeu de données} & \multicolumn{1}{c|}{\textbf{scRAW}} & \multicolumn{5}{c|}{\textbf{Rare Specific}} & \multicolumn{4}{c|}{\textbf{Généralistes}} & \multicolumn{4}{c}{\textbf{Correction Batch}} \\
    & \textbf{scRAW} & scAIDE & CellSIUS & DeepScena & scCAD & GiniClust & scVI & scMAE & PCA+Leiden & scNAME & Harmony & ComBat & DESC & Scanorama \\
    \midrule
    Zeisel & 0,611 & \underline{0,811} & 0,486 & 0,735 & 0,281 & 0,124 & 0,627 & 0,773 & 0,486 & \textbf{0,876} & 0,400 & 0,000 & 0,497 & 0,486 \\
    Spleen & 0,529 & 0,682 & 0,694 & 0,916 & 0,265 & 0,079 & 0,933 & 0,709 & \textbf{0,992} & 0,913 & \underline{0,990} & 0,978 & 0,933 & 0,799 \\
    Baron Pancreas & \underline{0,896} & 0,635 & 0,785 & 0,606 & 0,349 & 0,076 & \textbf{0,918} & 0,616 & 0,812 & 0,893 & 0,811 & 0,812 & 0,776 & 0,738 \\
    Human Testis & \textbf{0,967} & 0,844 & 0,813 & 0,358 & 0,291 & 0,110 & 0,915 & 0,847 & \underline{0,921} & 0,835 & 0,906 & 0,613 & 0,865 & 0,599 \\
    Kang PBMC & \textbf{0,678} & 0,220 & 0,194 & 0,263 & 0,220 & 0,006 & 0,158 & 0,216 & \underline{0,331} & 0,003 & 0,328 & 0,328 & 0,289 & 0,178 \\
    Macaque Retina & \underline{0,958} & 0,922 & 0,925 & 0,026 & 0,650 & 0,038 & 0,918 & 0,915 & 0,925 & 0,651 & 0,925 & \textbf{0,969} & 0,925 & 0,626 \\
    Bone Marrow & 0,126 & 0,364 & 0,329 & 0,111 & 0,279 & 0,009 & \textbf{0,499} & \underline{0,434} & 0,227 & 0,425 & 0,227 & 0,320 & 0,131 & 0,227 \\
    TM Liver & 0,323 & 0,135 & \textbf{0,543} & 0,027 & \underline{0,408} & 0,081 & 0,130 & 0,130 & 0,135 & 0,269 & 0,135 & 0,135 & 0,000 & 0,135 \\
    \midrule
    \textbf{Moyenne} & \underline{0,636} & 0,577 & 0,596 & 0,380 & 0,343 & 0,065 & \textbf{0,637} & 0,580 & 0,604 & 0,608 & 0,590 & 0,519 & 0,552 & 0,473 \\
    \bottomrule
  \end{tabular}
  }
  \caption{Résultats bruts de Précision sur Classes Rares (RareACC) pour les méthodes principales sur les 8 jeux de données communs. La meilleure performance par jeu de données est en gras, et la deuxième est soulignée.}
  \label{tab:common8_primary_raw_RareACC}
\end{table}

\begin{table}[H]
  \centering
  \fontsize{5.5pt}{7.0pt}\selectfont
  \setlength{\tabcolsep}{1.0pt}
  \begin{tabular}{l|c|ccccc|cccc|cccc}
    \toprule
    \textbf{Jeu de données} & \multicolumn{1}{c|}{\textbf{scRAW}} & \multicolumn{5}{c|}{\textbf{Rare Specific}} & \multicolumn{4}{c|}{\textbf{Généralistes}} & \multicolumn{4}{c}{\textbf{Correction Batch}} \\
    & \textbf{scRAW} & scAIDE & CellSIUS & DeepScena & scCAD & GiniClust & scVI & scMAE & PCA+Leiden & scNAME & Harmony & ComBat & DESC & Scanorama \\
    \midrule
    Zeisel & 0,556 & \underline{0,667} & 0,593 & 0,000 & 0,593 & 0,037 & 0,556 & 0,556 & 0,556 & \textbf{0,704} & 0,556 & 0,000 & 0,556 & 0,556 \\
    Spleen & \underline{0,976} & 0,000 & \textbf{1,000} & 0,000 & \textbf{1,000} & 0,024 & \textbf{1,000} & 0,000 & \underline{0,976} & 0,786 & \underline{0,976} & \underline{0,976} & 0,000 & \textbf{1,000} \\
    Baron Pancreas & \underline{0,814} & 0,627 & 0,636 & 0,000 & \textbf{0,915} & 0,576 & 0,466 & 0,695 & 0,797 & 0,432 & 0,788 & 0,788 & 0,466 & 0,754 \\
    Human Testis & 0,643 & \underline{0,929} & 0,607 & \textbf{0,964} & 0,571 & 0,000 & 0,000 & 0,000 & 0,000 & 0,000 & \underline{0,929} & 0,000 & 0,000 & 0,857 \\
    Kang PBMC & \underline{0,817} & 0,058 & 0,596 & -- & 0,274 & 0,000 & 0,037 & 0,043 & 0,562 & 0,000 & 0,540 & 0,544 & \textbf{0,870} & 0,093 \\
    Macaque Retina & 0,680 & 0,007 & 0,020 & 0,095 & 0,667 & 0,000 & 0,000 & 0,000 & 0,034 & 0,007 & 0,007 & \textbf{0,844} & 0,000 & \underline{0,762} \\
    Bone Marrow & 0,000 & 0,226 & 0,581 & 0,032 & \underline{0,677} & 0,000 & 0,000 & \textbf{0,742} & 0,000 & 0,032 & 0,000 & 0,000 & 0,000 & 0,000 \\
    TM Liver & -- & -- & -- & -- & -- & -- & -- & -- & -- & -- & -- & -- & -- & -- \\
    \midrule
    \textbf{Moyenne} & \underline{0,641} & 0,359 & 0,576 & 0,182 & \textbf{0,671} & 0,091 & 0,294 & 0,291 & 0,418 & 0,280 & 0,542 & 0,450 & 0,270 & 0,575 \\
    \bottomrule
  \end{tabular}
  \caption{Résultats bruts de Précision sur Classes Ultra-Rares (UltraRareACC) pour les méthodes principales sur les 8 jeux de données communs. La meilleure performance par jeu de données est en gras, et la deuxième est soulignée.}
  \label{tab:common8_primary_raw_UltraRareACC}
\end{table}

\begin{table}[H]
  \centering
  \fontsize{5.5pt}{7.0pt}\selectfont
  \setlength{\tabcolsep}{1.0pt}
  \resizebox{\textwidth}{!}{%
  \begin{tabular}{l|c|ccccc|cccc|cccc}
    \toprule
    \textbf{Jeu de données} & \multicolumn{1}{c|}{\textbf{scRAW}} & \multicolumn{5}{c|}{\textbf{Rare Specific}} & \multicolumn{4}{c|}{\textbf{Généralistes}} & \multicolumn{4}{c}{\textbf{Correction Batch}} \\
    & \textbf{scRAW} & scAIDE & CellSIUS & DeepScena & scCAD & GiniClust & scVI & scMAE & PCA+Leiden & scNAME & Harmony & ComBat & DESC & Scanorama \\
    \midrule
    Zeisel & 0,400 & 0,356 & 0,379 & \textbf{0,448} & 0,379 & 0,379 & \underline{0,442} & 0,408 & 0,379 & 0,382 & 0,254 & 0,058 & 0,292 & 0,378 \\
    Spleen & 0,791 & 0,690 & 0,564 & 0,645 & 0,565 & 0,556 & \textbf{0,837} & 0,551 & 0,554 & 0,675 & \underline{0,811} & 0,778 & 0,542 & 0,514 \\
    Baron Pancreas & 0,575 & 0,614 & 0,550 & 0,602 & 0,549 & 0,550 & 0,630 & 0,498 & 0,548 & 0,493 & \textbf{0,737} & 0,566 & \underline{0,697} & 0,530 \\
    Human Testis & 0,730 & \underline{0,782} & 0,667 & 0,775 & 0,667 & 0,667 & 0,775 & 0,699 & 0,662 & 0,679 & \textbf{0,797} & 0,562 & 0,778 & 0,634 \\
    Kang PBMC & 0,684 & 0,501 & 0,523 & 0,504 & 0,523 & 0,526 & 0,662 & 0,561 & \underline{0,726} & 0,545 & \underline{0,726} & \underline{0,726} & \textbf{0,795} & \underline{0,726} \\
    Macaque Retina & 0,637 & 0,451 & 0,508 & 0,517 & 0,508 & 0,513 & \textbf{0,683} & 0,382 & 0,510 & 0,366 & \underline{0,677} & 0,596 & 0,666 & 0,344 \\
    Bone Marrow & NA & NA & NA & NA & NA & NA & NA & NA & NA & NA & NA & NA & NA & NA \\
    TM Liver & NA & NA & NA & NA & NA & NA & NA & NA & NA & NA & NA & NA & NA & NA \\
    \midrule
    \textbf{Moyenne} & 0,636 & 0,566 & 0,532 & 0,582 & 0,532 & 0,532 & \textbf{0,671} & 0,517 & 0,563 & 0,523 & \underline{0,667} & 0,548 & 0,628 & 0,521 \\
    \bottomrule
  \end{tabular}
  }
  \caption{Résultats bruts de Correction de Batch pour les méthodes principales sur les 8 jeux de données communs. La meilleure performance par jeu de données est en gras, et la deuxième est soulignée. Les jeux sans effet de batch exploitable sont indiqués NA et exclus de la moyenne.}
  \label{tab:common8_primary_raw_Batch_correction}
\end{table}


\medskip
\noindent Pour faciliter la lecture des résultats, chaque sous-panneau de la figure~\ref{fig:scraw_common8_family_top3_plus_scraw} affiche \textbf{scRAW} ainsi que les trois meilleures méthodes de la famille pour la métrique considérée (basé sur la moyenne des huit datasets). Les méthodes restantes, ainsi que des résultats complémentaires avec méthode de correction de batch associées aux différents algorithmes, sont fournis dans cette annexe. Nous ne présentons pas dans le panneau principal les méthodes couplées à un outil de correction de batch, car les auteurs ne décrivent pas cet usage, ni comme prétraitement recommandé, ni comme composant du modèle. Ces variantes sont donc conservées uniquement comme analyses complémentaires. Pour les méthodes de deep learning généralistes, Harmony a été appliqué sur l'espace latent appris, avant le clustering final, car il s'agissait de la correction de batch la plus robuste dans les expériences de référence \cite{tran2020benchmark}. Cette pratique est justifiée par la conception même de Harmony : l'algorithme opère sur tout espace de faible dimension et ne suppose pas que celui-ci provienne d'une PCA obligatoirement. En revanche, aucun des papiers originaux des méthodes benchmarkées ici (scDeepCluster, scMAE, scNAME, scCDCG) ne propose ni n'évalue ce couplage de manière systématique : notre protocole comble ce manque en évaluant l'apport de Harmony sur l'espace latent de plusieurs architectures profondes avec des métriques globales et centrées sur les populations rares.

\begin{figure}[H]
    \centering
    \includegraphics[width=\textwidth,height=0.93\textheight,keepaspectratio]{scraw_common8_family_panel/common8_harmony_complement_panel.pdf}
    \caption{Résultats complémentaires avec Harmony. Les méthodes \textit{Rare Specific}+Harmony sont indiquées à titre exploratoire, car cette association n'est pas documentée comme usage standard par les auteurs. Pour les méthodes de deep learning traditionnelles, Harmony est appliqué sur l'espace latent avant le clustering final. La ligne Harmony correspond au pipeline PCA+Harmony avant clustering ; scVI est conservé comme méthode native intégrant l'information de batch.}
    \label{fig:common8_harmony_complement}
\end{figure}

Les tableaux suivants présentent les résultats bruts détaillés par jeu de données pour chaque métrique d'évaluation pour les expériences complémentaires avec Harmony :
\begin{table}[H]
  \centering
  \scriptsize
  \setlength{\tabcolsep}{1.5pt}
  \resizebox{\textwidth}{!}{%
  \begin{tabular}{l|c|ccccc|cccc}
    \toprule
    \textbf{Jeu de données} & \multicolumn{1}{c|}{\textbf{scRAW}} & \multicolumn{5}{c|}{\textbf{Rare + Harmony}} & \multicolumn{4}{c}{\textbf{Généralistes + Harmony}} \\
    & \textbf{scRAW} & scAIDE & CellSIUS & scCAD & DeepScena & GiniClust & Harmony & scVI & scMAE & scNAME \\
    \midrule
    Zeisel & 0,430 & 0,394 & 0,440 & 0,358 & 0,117 & 0,036 & 0,388 & \underline{0,513} & 0,129 & \textbf{0,722} \\
    Spleen & 0,870 & \textbf{0,931} & 0,558 & 0,416 & 0,487 & 0,141 & 0,518 & 0,572 & \underline{0,911} & 0,621 \\
    Baron Pancreas & \textbf{0,916} & \underline{0,914} & 0,895 & 0,605 & 0,440 & 0,052 & 0,851 & 0,621 & 0,770 & 0,728 \\
    Human Testis & 0,611 & 0,599 & 0,603 & 0,399 & 0,523 & 0,043 & \underline{0,652} & 0,594 & 0,651 & \textbf{0,653} \\
    Kang PBMC & 0,546 & \textbf{0,816} & 0,745 & 0,428 & 0,742 & 0,052 & 0,774 & \underline{0,805} & 0,760 & 0,492 \\
    Macaque Retina & 0,924 & \textbf{0,974} & 0,943 & 0,425 & 0,518 & 0,023 & 0,941 & \underline{0,965} & 0,950 & 0,498 \\
    Bone Marrow & 0,283 & 0,336 & \underline{0,341} & 0,224 & 0,222 & 0,007 & 0,312 & 0,283 & \textbf{0,386} & 0,325 \\
    TM Liver & \textbf{0,683} & \underline{0,647} & 0,365 & 0,443 & 0,566 & 0,103 & 0,535 & 0,519 & 0,491 & 0,477 \\
    \midrule
    \textbf{Moyenne} & \underline{0,658} & \textbf{0,701} & 0,611 & 0,412 & 0,452 & 0,057 & 0,621 & 0,609 & 0,631 & 0,565 \\
    \bottomrule
  \end{tabular}
  }
  \caption[Scores ARI avec Harmony (expériences complémentaires)]{Résultats bruts complémentaires de Rand Index Ajusté (ARI) avec Harmony pour les 8 jeux de données communs. La meilleure performance par jeu de données est en gras, et la deuxième est soulignée.}
  \label{tab:common8_harmony_complement_raw_ARI}
\end{table}

\begin{table}[H]
  \centering
  \scriptsize
  \setlength{\tabcolsep}{1.5pt}
  \begin{tabular}{l|c|ccccc|cccc}
    \toprule
    \textbf{Jeu de données} & \multicolumn{1}{c|}{\textbf{scRAW}} & \multicolumn{5}{c|}{\textbf{Rare + Harmony}} & \multicolumn{4}{c}{\textbf{Trad. + Harmony}} \\
    & \textbf{scRAW} & scAIDE & CellSIUS & scCAD & DeepScena & GiniClust & Harmony & scVI & scMAE & scNAME \\
    \midrule
    Zeisel & 0,641 & 0,625 & 0,523 & 0,343 & 0,363 & 0,297 & 0,606 & \underline{0,716} & 0,363 & \textbf{0,743} \\
    Spleen & 0,829 & \textbf{0,959} & 0,763 & 0,557 & 0,679 & 0,673 & 0,687 & 0,796 & \underline{0,955} & 0,822 \\
    Baron Pancreas & \textbf{0,934} & 0,874 & \underline{0,922} & 0,537 & 0,552 & 0,310 & 0,891 & 0,712 & 0,788 & 0,738 \\
    Human Testis & 0,718 & 0,705 & 0,669 & 0,315 & 0,595 & 0,217 & \underline{0,730} & 0,725 & 0,724 & \textbf{0,796} \\
    Kang PBMC & 0,706 & \textbf{0,833} & 0,731 & 0,408 & 0,772 & 0,421 & 0,753 & \underline{0,814} & 0,748 & 0,490 \\
    Macaque Retina & 0,937 & \textbf{0,984} & 0,967 & 0,560 & 0,570 & 0,220 & 0,965 & \underline{0,973} & 0,973 & 0,553 \\
    Bone Marrow & 0,416 & 0,476 & \underline{0,488} & 0,287 & 0,334 & 0,153 & 0,472 & 0,403 & \textbf{0,511} & 0,464 \\
    TM Liver & \textbf{0,746} & \underline{0,720} & 0,497 & 0,401 & 0,631 & 0,387 & 0,709 & 0,688 & 0,603 & 0,602 \\
    \midrule
    \textbf{Moyenne} & \underline{0,741} & \textbf{0,772} & 0,695 & 0,426 & 0,562 & 0,335 & 0,727 & 0,728 & 0,708 & 0,651 \\
    \bottomrule
  \end{tabular}
  \caption{Résultats bruts complémentaires de Précision Globale (ACC) avec Harmony pour les 8 jeux de données communs. La meilleure performance par jeu de données est en gras, et la deuxième est soulignée.}
  \label{tab:common8_harmony_complement_raw_ACC}
\end{table}

\begin{table}[H]
  \centering
  \scriptsize
  \setlength{\tabcolsep}{1.5pt}
  \resizebox{\textwidth}{!}{%
  \begin{tabular}{l|c|ccccc|cccc}
    \toprule
    \textbf{Jeu de données} & \multicolumn{1}{c|}{\textbf{scRAW}} & \multicolumn{5}{c|}{\textbf{Rare + Harmony}} & \multicolumn{4}{c}{\textbf{Généralistes + Harmony}} \\
    & \textbf{scRAW} & scAIDE & CellSIUS & scCAD & DeepScena & GiniClust & Harmony & scVI & scMAE & scNAME \\
    \midrule
    Zeisel & 0,611 & \textbf{0,649} & 0,454 & 0,281 & 0,324 & 0,124 & 0,400 & \underline{0,627} & 0,319 & 0,032 \\
    Spleen & 0,529 & 0,734 & 0,562 & 0,265 & 0,920 & 0,079 & \textbf{0,990} & \underline{0,933} & 0,721 & 0,910 \\
    Baron Pancreas & 0,896 & 0,688 & 0,880 & 0,349 & 0,674 & 0,076 & 0,811 & \underline{0,918} & \textbf{0,928} & 0,561 \\
    Human Testis & \textbf{0,967} & 0,814 & 0,498 & 0,291 & 0,489 & 0,110 & 0,906 & \underline{0,915} & 0,841 & 0,827 \\
    Kang PBMC & \textbf{0,678} & 0,441 & 0,438 & 0,220 & \underline{0,488} & 0,006 & 0,328 & 0,158 & 0,164 & 0,247 \\
    Macaque Retina & \textbf{0,958} & \underline{0,955} & 0,929 & 0,650 & 0,352 & 0,038 & 0,925 & 0,918 & 0,913 & 0,438 \\
    Bone Marrow & 0,126 & 0,379 & 0,329 & 0,279 & 0,209 & 0,009 & 0,227 & \textbf{0,499} & 0,434 & \underline{0,458} \\
    TM Liver & 0,323 & 0,135 & \textbf{0,543} & \underline{0,408} & 0,291 & 0,081 & 0,135 & 0,130 & 0,130 & 0,004 \\
    \midrule
    \textbf{Moyenne} & \underline{0,636} & 0,599 & 0,579 & 0,343 & 0,468 & 0,065 & 0,590 & \textbf{0,637} & 0,556 & 0,435 \\
    \bottomrule
  \end{tabular}
  }
  \caption[Scores RareACC avec Harmony (expériences complémentaires)]{Résultats bruts complémentaires de Précision sur Classes Rares (RareACC) avec Harmony pour les 8 jeux de données communs. La meilleure performance par jeu de données est en gras, et la deuxième est soulignée.}
  \label{tab:common8_harmony_complement_raw_RareACC}
\end{table}

\begin{table}[H]
  \centering
  \scriptsize
  \setlength{\tabcolsep}{1.5pt}
  \resizebox{\textwidth}{!}{%
  \begin{tabular}{l|c|ccccc|cccc}
    \toprule
    \textbf{Jeu de données} & \multicolumn{1}{c|}{\textbf{scRAW}} & \multicolumn{5}{c|}{\textbf{Rare + Harmony}} & \multicolumn{4}{c}{\textbf{Généralistes + Harmony}} \\
    & \textbf{scRAW} & scAIDE & CellSIUS & scCAD & DeepScena & GiniClust & Harmony & scVI & scMAE & scNAME \\
    \midrule
    Zeisel & \underline{0,556} & 0,519 & 0,296 & \textbf{0,593} & 0,370 & 0,037 & \underline{0,556} & \underline{0,556} & 0,519 & 0,185 \\
    Spleen & \underline{0,976} & 0,000 & \textbf{1,000} & \textbf{1,000} & 0,286 & 0,024 & \underline{0,976} & \textbf{1,000} & 0,000 & 0,833 \\
    Baron Pancreas & \underline{0,814} & 0,780 & 0,627 & \textbf{0,915} & 0,466 & 0,576 & 0,788 & 0,466 & 0,678 & 0,458 \\
    Human Testis & 0,643 & \textbf{0,929} & \underline{0,893} & 0,571 & \underline{0,893} & 0,000 & \textbf{0,929} & 0,000 & 0,000 & 0,000 \\
    Kang PBMC & \textbf{0,817} & 0,058 & \underline{0,788} & 0,274 & -- & 0,000 & 0,540 & 0,037 & 0,038 & 0,308 \\
    Macaque Retina & \underline{0,680} & \textbf{0,878} & 0,007 & 0,667 & 0,000 & 0,000 & 0,007 & 0,000 & 0,000 & 0,000 \\
    Bone Marrow & 0,000 & 0,194 & 0,581 & \underline{0,677} & 0,097 & 0,000 & 0,000 & 0,000 & \textbf{0,742} & \textbf{0,742} \\
    TM Liver & -- & -- & -- & -- & -- & -- & -- & -- & -- & -- \\
    \midrule
    \textbf{Moyenne} & \underline{0,641} & 0,479 & 0,599 & \textbf{0,671} & 0,352 & 0,091 & 0,542 & 0,294 & 0,282 & 0,361 \\
    \bottomrule
  \end{tabular}
  }
  \caption[Scores UltraRareACC avec Harmony (expériences complémentaires)]{Résultats bruts complémentaires de Précision sur Classes Ultra-Rares (UltraRareACC) avec Harmony pour les 8 jeux de données communs. La meilleure performance par jeu de données est en gras, et la deuxième est soulignée.}
  \label{tab:common8_harmony_complement_raw_UltraRareACC}
\end{table}

\begin{table}[H]
  \centering
  \scriptsize
  \setlength{\tabcolsep}{1.5pt}
  \resizebox{\textwidth}{!}{%
  \begin{tabular}{l|c|ccccc|cccc}
    \toprule
    \textbf{Jeu de données} & \multicolumn{1}{c|}{\textbf{scRAW}} & \multicolumn{5}{c|}{\textbf{Rare + Harmony}} & \multicolumn{4}{c}{\textbf{Généralistes + Harmony}} \\
    & \textbf{scRAW} & scAIDE & CellSIUS & scCAD & DeepScena & GiniClust & Harmony & scVI & scMAE & scNAME \\
    \midrule
    Zeisel & 0,400 & 0,312 & 0,270 & 0,270 & 0,329 & 0,270 & 0,254 & \textbf{0,442} & 0,388 & \underline{0,402} \\
    Spleen & 0,791 & 0,756 & 0,785 & 0,783 & \underline{0,822} & 0,790 & 0,811 & \textbf{0,837} & 0,767 & 0,310 \\
    Baron Pancreas & 0,575 & 0,562 & 0,730 & 0,729 & 0,681 & \underline{0,730} & \textbf{0,737} & 0,630 & 0,604 & 0,309 \\
    Human Testis & 0,730 & 0,723 & 0,776 & 0,777 & 0,712 & \underline{0,777} & \textbf{0,797} & 0,775 & 0,752 & 0,430 \\
    Kang PBMC & \underline{0,684} & 0,573 & 0,664 & 0,663 & 0,602 & 0,664 & \textbf{0,726} & 0,662 & 0,561 & 0,423 \\
    Macaque Retina & 0,637 & 0,540 & 0,668 & 0,669 & 0,600 & 0,666 & \underline{0,677} & \textbf{0,683} & 0,495 & 0,229 \\
    Bone Marrow & NA & NA & NA & NA & NA & NA & NA & NA & NA & NA \\
    TM Liver & NA & NA & NA & NA & NA & NA & NA & NA & NA & NA \\
    \midrule
    \textbf{Moyenne} & 0,636 & 0,578 & 0,649 & 0,649 & 0,624 & 0,650 & \underline{0,667} & \textbf{0,671} & 0,594 & 0,350 \\
    \bottomrule
  \end{tabular}
  }
  \caption[Scores de correction de batch avec Harmony (expériences complémentaires)]{Résultats bruts complémentaires de Correction de Batch avec Harmony pour les 8 jeux de données communs. La meilleure performance par jeu de données est en gras, et la deuxième est soulignée. Les jeux sans effet de batch exploitable sont indiqués NA et exclus de la moyenne.}
  \label{tab:common8_harmony_complement_raw_Batch_correction}
\end{table}


Ces résultats détaillés permettent d'analyser le comportement individuel de chaque méthode sur chacun des jeux de données, offrant ainsi une vision plus fine et contrastée que les seules moyennes globales.

\newpage
\section{Étude d'ablation de scRAW}
\label{annexe:ablation_scraw}

Cette section présente l'étude d'ablation menée pour valider l'apport individuel de chaque composante de \textbf{scRAW} (la pondération dynamique, la perte triplet rare-aware et la correction adversariale de batch DANN). Toutes les variantes ont été entraînées sur le jeu de données Baron Human Pancreas en utilisant HDBSCAN comme algorithme de clustering final.

Le Tableau~\ref{tab:scraw_ablation} détaille les scores obtenus pour les variantes suivantes :
\begin{itemize}
    \item \textbf{MSE} : Autoencodeur classique avec une perte MSE standard sans pondération, sans triplet loss et sans correction adversariale.
    \item \textbf{Weighted MSE} : Variante incluant la pondération dynamique des cellules par taille de pseudo-cluster et densité locale, mais sans triplet loss ni correction adversariale.
    \item \textbf{Weighted + triplet} : Autoencodeur pondéré avec triplet loss rare-aware, sans correction de batch adversariale.
    \item \textbf{Weighted + DANN} : Autoencodeur pondéré avec correction adversariale (DANN), sans triplet loss.
    \item \textbf{Full (scRAW)} : Modèle complet associant la pondération dynamique, la triplet loss et la branche adversariale.
    \item \textbf{Full single update} : Modèle scRAW complet, mais où la mise à jour des poids cellulaires n'est effectuée qu'une seule fois au lieu d'être actualisée périodiquement.
    \item \textbf{Full density-only} : Modèle scRAW complet, mais où la pondération des cellules repose uniquement sur la densité locale (la composante liée à la taille des pseudo-clusters est désactivée).
    \item \textbf{Full cluster-size-only} : Modèle scRAW complet, mais où la pondération repose uniquement sur la taille des pseudo-clusters Leiden (la composante liée à la densité locale est désactivée).
\end{itemize}

\begin{table}[H]
\centering
\small
\renewcommand{\arraystretch}{1.15}
\begin{tabularx}{\textwidth}{X c c c c c c c r}
\toprule
\textbf{Variante} & \textbf{ARI} & \textbf{NMI} & \textbf{ACC} & \textbf{B. ACC} & \textbf{Rare} & \textbf{U. Rare} & \textbf{scIB-E} & \textbf{Temps (s)} \\
\midrule
MSE & 0,189 & 0,397 & 0,407 & 0,427 & 0,350 & 0,695 & 0,427 & 70,8 \\
Weighted MSE & 0,269 & 0,455 & 0,456 & 0,520 & 0,390 & 0,686 & 0,426 & 518,2 \\
Weighted + triplet & 0,851 & 0,871 & 0,885 & 0,813 & 0,809 & 0,805 & 0,431 & 454,9 \\
Weighted + DANN & 0,246 & 0,560 & 0,545 & 0,611 & 0,491 & 0,653 & 0,574 & 279,7 \\
\midrule
\textbf{Full (scRAW)} & \textbf{0,916} & \textbf{0,884} & \textbf{0,934} & \textbf{0,870} & \textbf{0,896} & \textbf{0,814} & \textbf{0,590} & \textbf{337,9} \\
Full single update & 0,900 & 0,873 & 0,913 & 0,747 & 0,793 & 0,686 & 0,599 & 285,9 \\
Full density-only & 0,920 & 0,888 & 0,926 & 0,753 & 0,796 & 0,703 & 0,587 & 334,4 \\
Full cluster-size-only & 0,723 & 0,777 & 0,800 & 0,869 & 0,897 & 0,949 & 0,596 & 347,2 \\
\bottomrule
\end{tabularx}
\caption{Résultats de l'étude d'ablation des composantes de scRAW sur le jeu de données Baron Human Pancreas (clustering final HDBSCAN). B. ACC : BalancedACC ; Rare : RareACC ; U. Rare : UltraRareACC ; scIB-E : score d'intégration de batch scIB (entropie de mélange).}
\label{tab:scraw_ablation}
\end{table}

\begin{figure}[H]
    \centering
    \includegraphics[width=0.9\textwidth]{stable_default_mse_ablations_20260526_210436/figures/stable_default_mse_ablation_key_metrics_barplot.png}
    \caption{Comparaison graphique des performances des variantes d'ablation de scRAW pour les métriques de clustering global (ARI, NMI, ACC), la détection des classes rares (RareACC, UltraRareACC) et le score d'intégration de batch (scIB-E).}
    \label{fig:scraw_ablation_barplot}
\end{figure}

\subsection*{Discussion des résultats}

L'analyse de ces résultats d'ablation permet de tirer plusieurs enseignements essentiels sur l'architecture de scRAW :
\begin{itemize}
    \item \textbf{Importance de la pondération de la reconstruction} : L'autoencodeur classique (\textit{MSE}) obtient des scores très limités (ARI $= 0{,}189$ et NMI $= 0{,}397$). L'activation de la pondération dynamique (\textit{Weighted MSE}) augmente nettement ces scores (ARI $= 0{,}269$ et NMI $= 0{,}455$) ainsi que la BalancedACC ($0{,}520$ contre $0{,}427$), ce qui confirme que moduler l'importance des cellules selon leur rareté et leur densité locale aide à structurer l'espace latent de manière plus équitable.
    \item \textbf{Rôle fondamental de la triplet loss rare-aware} : L'ajout de la perte triplet (\textit{Weighted + triplet}) engendre une amélioration spectaculaire sur l'ensemble des critères. L'ARI bondit à $0{,}851$, la BalancedACC à $0{,}813$, et les exactitudes sur les populations rares et ultra-rares dépassent $0{,}80$. Cela montre que guider explicitement l'organisation géométrique des embeddings en rapprochant les cellules similaires et en isolant les petits groupes est indispensable pour obtenir une partition biologique claire et robuste.
    \item \textbf{Synergie avec la correction adversariale} : L'utilisation de DANN seul avec la pondération (\textit{Weighted + DANN}) améliore considérablement le score de correction de batch (l'entropie scIB-E grimpe à $0{,}574$) mais reste inefficace pour former des groupes biologiquement cohérents sans le terme de triplet loss. Le modèle complet (\textit{Full}) réunit le meilleur des deux aspects : il maintient une performance de clustering très élevée (ARI $= 0{,}916$, ACC $= 0{,}934$) tout en garantissant un excellent niveau d'intégration de batch (scIB-E $= 0{,}590$).
    \item \textbf{Rythme de mise à jour des poids} : Dans la variante \textit{Full single update}, recalculer les poids une seule fois plutôt que périodiquement dégrade la détection des populations les plus petites (l'Ultra-rare ACC chute de $0{,}814$ à $0{,}686$). L'actualisation continue des poids au fil de l'apprentissage est donc cruciale pour s'adapter à l'évolution de la topologie de l'espace latent.
    \item \textbf{Complémentarité des composantes de pondération} : Se restreindre à la seule densité locale (\textit{Full density-only}) offre la meilleure ARI globale ($0{,}920$) mais diminue la préservation des classes ultra-rares (Ultra-rare ACC $= 0{,}703$). Inversement, la pondération par taille de pseudo-cluster seule (\textit{Full cluster-size-only}) favorise massivement les populations très petites (Ultra-rare ACC $= 0{,}949$) mais dégrade le partitionnement global (ARI $= 0{,}723$). La combinaison multiplicative des deux termes (modèle complet \textit{Full}) permet d'atteindre le compromis le plus équilibré et le plus robuste.
\end{itemize}

\newpage
\section{Hyperparamètres de la configuration de référence scRAW}
\label{annexe:hyperparametres_scraw}

Cette annexe détaille les hyperparamètres effectivement actifs dans la configuration par défaut utilisée pour décrire scRAW dans le rapport.

{\scriptsize
\setlength{\tabcolsep}{3pt}
\renewcommand{\arraystretch}{1.15}
\begin{longtable}{p{0.15\textwidth}p{0.27\textwidth}p{0.20\textwidth}p{0.30\textwidth}}
\caption{Hyperparamètres actifs du preset \texttt{stable\_default} de scRAW.}
\label{tab:scraw_stable_default_hyperparams}\\
\toprule
\textbf{Bloc} & \textbf{Paramètre} & \textbf{Valeur} & \textbf{Rôle} \\
\midrule
\endfirsthead
\toprule
\textbf{Bloc} & \textbf{Paramètre} & \textbf{Valeur} & \textbf{Rôle} \\
\midrule
\endhead
\midrule
\multicolumn{4}{r}{\textit{Suite page suivante}}\\
\endfoot
\bottomrule
\endlastfoot

Prétraitement & \texttt{n\_top\_genes} & 2000 & Nombre de gènes hautement variables conservés. \\
Prétraitement & \texttt{min\_genes\_per\_cell} & 200 & Filtrage des cellules de mauvaise qualité. \\
Prétraitement & \texttt{min\_cells\_per\_gene} & 3 & Filtrage des gènes très peu détectés. \\
Prétraitement & \texttt{target\_sum} & 20000 & Normalisation des comptes par cellule. \\
Prétraitement & \texttt{scale\_max\_value} & 10.0 & Seuil appliqué après standardisation. \\
Prétraitement & \texttt{hvg\_flavor} & \texttt{seurat} & Méthode de sélection des HVG. \\

Architecture & \texttt{hidden\_layers} & \texttt{512,256,128} & Tailles des couches cachées de l'encodeur. \\
Architecture & \texttt{z\_dim} & 256 & Dimension de l'espace latent. \\
Architecture & \texttt{dropout} & 0.3 & Dropout dans l'encodeur et le décodeur. \\

Entraînement & \texttt{epochs} & 120 & Nombre total d'époques. \\
Entraînement & \texttt{warmup\_epochs} & 55 & Phase sans pondération dynamique. \\
Entraînement & \texttt{lr} & 0.00164076083297036 & Taux d'apprentissage initial. \\
Entraînement & \texttt{batch\_size} & 192 & Taille des mini-batchs. \\
Entraînement & \texttt{seed} & 64 & Graine utilisée pour l'entraînement. \\

Reconstruction & \texttt{reconstruction\_distribution} & \texttt{mse} & Perte de reconstruction utilisée. \\
Reconstruction & \texttt{masking\_rate} & 0.1 & Fraction d'entrées masquées aléatoirement. \\
Reconstruction & \texttt{masking\_value} & 0.0 & Valeur utilisée pour les entrées masquées. \\
Reconstruction & \texttt{masked\_recon\_weight} & 0.8 & Poids interne des positions masquées. \\
Reconstruction & \texttt{masking\_apply\_weighted} & \texttt{True} & Masquage aussi actif pendant la phase pondérée. \\

Pondération & \texttt{weight\_exponent} & 0.2 & Intensité de la pondération par taille de pseudo-cluster. \\
Pondération & \texttt{density\_knn\_k} & 15 & Voisinage utilisé pour estimer la densité locale. \\
Pondération & \texttt{density\_weight\_exponent} & 1.0 & Exposant de la composante densité. \\
Pondération & \texttt{density\_weight\_clip} & 3.0 & Borne supérieure avant normalisation de la densité. \\
Pondération & \texttt{weight\_fusion\_mode} & \texttt{multiplicative} & Fusion des composantes cluster et densité. \\
Pondération & \texttt{cluster\_weight\_power} & 1.0 & Puissance appliquée à la composante cluster. \\
Pondération & \texttt{density\_weight\_power} & 1.0 & Puissance appliquée à la composante densité. \\
Pondération & \texttt{dynamic\_weight\_momentum} & 0.6884621079434989 & Lissage des poids entre deux mises à jour. \\
Pondération & \texttt{dynamic\_weight\_update\_interval} & 20 & Fréquence de recalcul des poids. \\
Pondération & \texttt{min\_cell\_weight} & 0.3845423008053828 & Borne inférieure du poids cellulaire final. \\
Pondération & \texttt{max\_cell\_weight} & 10.0 & Borne supérieure du poids cellulaire final. \\

Pseudo-labels & \texttt{pseudo\_label\_method} & \texttt{leiden} & Méthode utilisée pour les pseudo-clusters. \\
Pseudo-labels & \texttt{pseudo\_leiden\_resolution\_strategy} & \texttt{scraw\_default} & Recherche de résolution Leiden utilisée par scRAW. \\
Pseudo-labels & \texttt{unsupervised\_k\_selection} & \texttt{heuristic} & Estimation heuristique de \(K\) si nécessaire. \\
Pseudo-labels & \texttt{unsupervised\_k\_min} & 8 & Borne minimale de \(K\). \\
Pseudo-labels & \texttt{unsupervised\_k\_max} & 30 & Borne maximale de \(K\). \\

Triplet & \texttt{rare\_triplet\_weight} & 0.05007581780188212 & Poids de la triplet loss. \\
Triplet & \texttt{rare\_triplet\_start\_epoch} & 60 & Début de la triplet loss. \\
Triplet & \texttt{rare\_triplet\_margin} & 0.4 & Marge de séparation. \\
Triplet & \texttt{rare\_triplet\_min\_weight} & 1.2 & Poids minimal pour choisir une ancre. \\
Triplet & \texttt{max\_triplet\_anchors\_per\_batch} & 64 & Nombre maximal d'ancres par mini-batch. \\

Batch & \texttt{use\_batch\_conditioning} & \texttt{True} & Activation de la branche batch si plusieurs lots existent. \\
Batch & \texttt{batch\_correction\_key} & \texttt{batch} & Colonne de métadonnées utilisée pour le batch. \\
Batch & \texttt{adversarial\_batch\_weight} & 0.11763398875166495 & Poids de la perte adversariale batch. \\
Batch & \texttt{adversarial\_lambda} & 1.0 & Intensité de l'inversion de gradient. \\
Batch & \texttt{adversarial\_start\_epoch} & 30 & Début de la branche adversariale. \\
Batch & \texttt{adversarial\_ramp\_epochs} & 30 & Durée de montée progressive du signal adversarial. \\

Clustering final & \texttt{hdbscan\_min\_cluster\_size} & 8 & Taille minimale d'un cluster HDBSCAN. \\
Clustering final & \texttt{hdbscan\_min\_samples} & 6 & Nombre minimal d'échantillons HDBSCAN. \\
Clustering final & \texttt{hdbscan\_cluster\_selection\_method} & \texttt{eom} & Stratégie de sélection des clusters. \\
Clustering final & \texttt{hdbscan\_reassign\_noise} & \texttt{False} & Pas de réassignation des points bruit. \\
\end{longtable}
}

\newpage
\section{Figures de recherche d'hyperparamètres scRAW}
\label{annexe:figures_hyperparametres_scraw}

Cette annexe regroupe les sorties Optuna d'importance moyenne des hyperparamètres pour les trois recherches utilisées dans la construction de \texttt{stable\_default}.

\begin{figure}[H]
    \centering
    \includegraphics[width=\textwidth,height=0.72\textheight,keepaspectratio]{hyperparametres_scraw/baron_param_importances_mean5_top30.png}
    \caption{Importance des hyperparamètres pour la recherche large sur Baron Human Pancreas.}
    \label{fig:scraw_hparam_importance_baron}
\end{figure}

\begin{figure}[H]
    \centering
    \includegraphics[width=\textwidth,height=0.72\textheight,keepaspectratio]{hyperparametres_scraw/generalist8_param_importances_mean5_top30.png}
    \caption{Importance des hyperparamètres pour la recherche généraliste sur huit jeux de données.}
    \label{fig:scraw_hparam_importance_generalist8}
\end{figure}

\begin{figure}[H]
    \centering
    \includegraphics[width=\textwidth,height=0.72\textheight,keepaspectratio]{hyperparametres_scraw/stable_generalist_param_importances_mean5_top30.png}
    \caption{Importance des hyperparamètres pour la recherche restreinte \texttt{stable\_generalist}.}
    \label{fig:scraw_hparam_importance_stable_generalist}
\end{figure}

\newpage
\section{Screenshots de l'interface graphique SCRBenchmark}
\label{annexe:screenshots}

Cette section regroupe les captures d'écran illustrant le fonctionnement et le flux complet de l'interface utilisateur Streamlit développée pour SCRBenchmark.

\begin{figure}[H]
    \centering\footnotesize
    \begin{subfigure}[b]{0.45\textwidth}
        \centering
        \includegraphics[width=\textwidth]{SCRBenchmark/scr_01_upload.png}
        \caption{Importation des données}
        \label{fig:scr_upload}
    \end{subfigure}
    \hfill
    \begin{subfigure}[b]{0.45\textwidth}
        \centering
        \includegraphics[width=\textwidth]{SCRBenchmark/scr_02_visualisation_des_donnees.png}
        \caption{Visualisation des données}
        \label{fig:scr_vis_donnees}
    \end{subfigure}
    
    \vspace{0.3cm}
    
    \begin{subfigure}[b]{0.45\textwidth}
        \centering
        \includegraphics[width=\textwidth]{SCRBenchmark/scr_03_choixExperience.png}
        \caption{Choix de l'expérience}
        \label{fig:scr_choix_exp}
    \end{subfigure}
    \hfill
    \begin{subfigure}[b]{0.45\textwidth}
        \centering
        \includegraphics[width=\textwidth]{SCRBenchmark/scr_04_preprocess.png}
        \caption{Prétraitement}
        \label{fig:scr_preprocess}
    \end{subfigure}
    
    \vspace{0.3cm}
    
    \begin{subfigure}[b]{0.45\textwidth}
        \centering
        \includegraphics[width=\textwidth]{SCRBenchmark/scr_05_choixAlgo.png}
        \caption{Choix des algorithmes}
        \label{fig:scr_choix_algo}
    \end{subfigure}
    \hfill
    \begin{subfigure}[b]{0.45\textwidth}
        \centering
        \includegraphics[width=\textwidth]{SCRBenchmark/scr_06_CommandeCLI.png}
        \caption{Génération CLI}
        \label{fig:scr_cli}
    \end{subfigure}
    
    \caption{Aperçu des différentes étapes de l'interface Streamlit de SCRBenchmark (1\textsuperscript{re} partie : configuration et préparation).}
    \label{fig:scrbenchmark_panel}
\end{figure}

\begin{figure}[H]
    \ContinuedFloat
    \centering\footnotesize
    \begin{subfigure}[b]{0.45\textwidth}
        \centering
        \includegraphics[width=\textwidth]{SCRBenchmark/scr_07_runner.png}
        \caption{Exécution des calculs}
        \label{fig:scr_runner}
    \end{subfigure}
    \hfill
    \begin{subfigure}[b]{0.45\textwidth}
        \centering
        \includegraphics[width=\textwidth]{SCRBenchmark/scr_08_results.png}
        \caption{Résultats des métriques}
        \label{fig:scr_results}
    \end{subfigure}
    
    \vspace{0.3cm}
    
    \begin{subfigure}[b]{0.45\textwidth}
        \centering
        \includegraphics[width=\textwidth]{SCRBenchmark/scr_09_resultsUmaps.png}
        \caption{UMAPs des résultats}
        \label{fig:scr_results_umaps}
    \end{subfigure}
    
    \caption{Aperçu des différentes étapes de l'interface Streamlit de SCRBenchmark (2\textsuperscript{e} partie : exécution et résultats).}
\end{figure}
